\documentclass[11pt,a4paper]{article}
\usepackage[utf8]{inputenc}
\usepackage[T1]{fontenc}
\usepackage{amsmath,amssymb,amsthm}
\usepackage{geometry}
\geometry{margin=2.4cm}

\newtheorem{definition}{Definition}
\newtheorem{remark}{Remark}

\title{Exploratory Coupling Ratios and Logarithmic Scales Against the Fine-Structure Constant $\alpha$ Under OVRP}
\author{Mariano Panzano Caball\'e \\ \small Oasis Sovereign Architecture Research, Barcelona}
\date{November 26, 2026}

\begin{document}
\maketitle

\begin{abstract}
Expanding on Zenodo draft \texttt{19339999} and connecting with the \emph{Oasis Verifiable Research Pipeline} (\texttt{DOI: 10.5281/zenodo.24443531}), logarithmic scale audits (\texttt{DOI: 10.5281/zenodo.24538777}), and ternary/Q-OS notes (\texttt{DOI: 10.5281/zenodo.24559695}, \texttt{DOI: 10.5281/zenodo.24560686}), this study evaluates candidate geometric-logarithmic expressions against the CODATA reference inverse fine-structure constant $1/\alpha_{\mathrm{CODATA}} \approx 137.035999$. All numerical residuals are reported explicitly alongside a zero-\texttt{sorry} \texttt{Lean~4.36.0} specification (\texttt{Paper28\_AlphaCoupling\_Spec.lean}), a \texttt{C99} probe, \texttt{SHA-256} and Bitcoin OpenTimestamps provenance, AI assistance disclosure, and explicit open boundaries under \texttt{CC-BY-4.0}.
\end{abstract}

\section{Cross-Linkage and Candidate Ratio Audit}
In standard quantum electrodynamics, $\alpha$ is empirically determined and no first-principles analytical derivation from mathematical constants is established. Within OVRP, candidate combinations of $\pi$, $\phi = (1+\sqrt{5})/2$, $\ln 2$, $\ln 3$, and $\ln 10$ are tested strictly as numerical probes and comparative benchmarks rather than proven physical laws.

\section*{AI Assistance Disclosure and Open Boundaries}
Prepared with AI assistance under author direction and verification, paired with machine checking in \texttt{Lean~4.36.0} (\texttt{0~sorry}) and a deterministic \texttt{C99} probe under \texttt{CC-BY-4.0}.

\begin{remark}[Explicit Open Boundaries]
Numerical proximity or geometric ratios involving $\pi$, $\phi$, $\ln 2$, $\ln 3$, or $\ln 10$ do not constitute a first-principles derivation of the fine-structure constant $\alpha$ or resolve open problems in fundamental physics.
\end{remark}

\end{document}
